% !TeX root = ../papers/nc-p3.tex
\section{有限点グラフの連結トーラスと二つの重み付き局所型}
\label{sec:finite-torus-profiles}
\begin{theorem}\label{thm:finite-torus-profiles}
四辺の点対応が有限な標準正則窓では、各辺の六関係を保つ射影群
$\PGL(V_i)\times\PGL(V_{i+1})$ の部分群のトーラス階数は高々1である。
窓全体を保つ非自明な乗法群があれば、スカラーと共通倍率を除き四位置の矢および関係の重みは
全て等しく、可能性は下表に限られる。
\[
\begin{array}{c|c|c|l}
 V & R & h\text{ の指数}&\text{点スキームの長さ}\\\hline
 (0,1,2,3)&(1,2,3,3,4,5)&6&\text{最低・最高の二点に各10}\\
 (0,1,1,2)&(1,1,2,2,3,3)&4&\text{最低・最高に各9、中間の固定直線上に計2}
\end{array}
\]
後者の中間部分は二つの射影直線の積で二本の $(1,1)$ 形式が切る全交叉である。
いずれも必要条件で、値の列だけで正則性は主張しない。
\end{theorem}

\begin{proof}[露出面の次元と二端の重み差]
各辺の六空間 $R$ は純テンソルを含まない（定理\ref{thm:finite-frobenius-low-opposite-as}）。
最小重みの矢部分の次元を二端で $a,b$ とする。最小和の関係空間はこの $a\times b$ ブロックに
そのまま含まれ、Segreを避けるので次元は $(a-1)(b-1)$ 以下である。
二つの固有空間の射影双対の積に制限するとこの関係だけが残る。
有限交叉のためには少なくとも $a+b-2$ 本が必要である。従って $a=b=1$ または $a,b\ge3$ が必要で、特に2は現れない。

辺を保つトーラスは各位置で忠実に作用する。片端の作用が恒等なら、有限グラフのもう片端に
誘導される作用も点スキーム上恒等である。よって $x\wedge gx$ の各二次式はその上で消え、
定理\ref{thm:finite-square-chart}の二次制限の単射性から $g$ はスカラーである。
有効階数が2以上なら重みの凸包は二次元以上である。頂点の重複度が2ならその頂点を露出させれば矛盾する。
それ以外は三角形・四角形・四面体の中に、重複込みでちょうど二つの重みを持つ辺がある
（三角形に残りの一点があっても高々一辺だけに入る）。
この辺を露出させる一径数を使っても同じ矛盾となる。

窓作用に戻る。全位置で最低重みを0にする。最低または最高が三重なら全位置で同じ側が三重であり、
定理\ref{thm:binary-torus-windows}の $S[z]$ になる。これは有限点グラフを持たない。
実際 $z=0$ に三生成商の点対応が残り、三本の $(1,1)$ 形式を二平面の積に課しても次元は1以上である。
従って両極は単純である。各辺の最低どうしは点対となり、そこでの各側の評価は階数3である。
二端の最初の正の重みが違えば、小さい方の重みで可能な関係は片方が最低直線である純ブロックだけに入り、
そこには非零関係がない。しかし最低点で相手の全ての正の重みが評価像に必要である。
従って最初の差 $d$ は等しい。その重みでの二ブロックは各々片側が直線であり、
純関係がないことから両射影は単射、最低点での評価から両射影は全射となる。
ゆえに第二重みの重複度も等しい。反転すれば最後の差 $e$ とその重複度も全位置で共通である。
二重みだけなら定理\ref{thm:binary-torus-windows}で既に尽き、$2+2$ も純軸の点直線を持つので有限でない。
\end{proof}

\begin{proof}[単純四重みと固定点の重複度]
隣接矢の重みを $a=(0,d,d+s,d+s+e)$, $b=(0,d,d+t,d+t+e)$ と置く。
各和 $a_i+b_j$ は $(d,e,s,t)$ の整数一次式である。実際に成立する和の一致を全て保つ
新しい正の四数は、各関係重み空間を細分しないので、同じ辺の別の作用を与える。
整数解空間から射影二位置への写像は単射で、前段の階数上界から等号の行列は階数3である。
正の解を許す18枚の異なる超平面から三枚ずつを選び、符号付き三次小行列式を取ると、
正の原始候補は31本である。

各候補では $c_z=\#\{(i,j):a_i+b_j=z\}$ として重み $z$ の関係数 $0\le r_z<c_z$ を使う
（存在する和だけを数える）。全ブロックを関係にすることは純関係を含むため不可能である。
$r_{a_i+b_j}=0$ なら座標点対 $(i,j)$ は固定点である。逆も、局所指標の次式から従う。
その点対における局所環の指標（点がなければ零）は
\[
 L_{ij}(u)=
 \frac{\prod_z(1-u^{z-a_i-b_j})^{r_z}}
 {\prod_{k\ne i}(1-u^{a_k-a_i})\prod_{l\ne j}(1-u^{b_l-b_j})}.
\]
実際 $(n,m)$ を十分大きくし、六本の完全交叉Koszul分解の指標を取る。
完全対称式 $\operatorname{ch}\Sym^n V=\sum_i u^{n a_i}/\prod_{k\ne i}(1-u^{a_k-a_i})$
を両位置に入れると $u^{na_i+mb_j}$ の係数がこの式となる。
他方、有限スキームの全閉点は固定され、各局所因子で座標自明化により同じ係数は局所環指標となる。
異なる $(i,j)$ の係数は $n,m$ を動かしてVandermondeで分離できる。
これで $r_{a_i+b_j}>0$ の対には点がない。残る点対は部分マッチングをなし、$L_{ij}(1)$ は正整数で総和20である。

\path{computations/verify_finite_torus_profiles.py} がこれらの条件を全31候補の全六関係で比較する。
結果は $(d,e)=(1,1)$ として $(s,t)=(1,1),(1,2),(2,1)$ だけで、いずれも両極の点対が各長さ10である。
最初の関係重みは表の上段、他の二つでは $1,2,3,4,5,6$ である。
全四位置で倍率が共通なのは $d$ が共通だからである。
中央差2の位置数 $j$ は隣接禁止から高々2、矢の行列式文字の等式から $4H=24+2j$ となる。
従って $j=2$ なら交互に現れ $H=7$。この場合のEuler第一べき和は重み2の係数が $-2$ で、
四文字の指標にならない。以上で上段が従う。
\end{proof}

\begin{proof}[中央二重の四重み]
全位置が $(0,d,d,d+e)$ で共通である。$d+e\ne2d$ と仮定すると、重み $2d$ の関係は
二つの中央二平面の積に純粋に入り、その次元はSegreを避けるため高々1である。
両端で最高文字を零にして最低座標を1と置く。重み $d$ の二関係は中央二変数間の可逆な線形対応を与え、
残る式は重み $2d$ の高々一つの斉次二次式のみとなる。二変数上なので正次元の解が生じ、有限性に反する。
従って原始化すれば $d=e=1$。重み1と3の各二関係に加え重み2が二つで、他の成分は零である。
最低点のアフィン座標の重みは両端とも $1,1,2$ で、関係の重みは $1,1,2,2,3,3$ である。
正に次数付けられた六変数上の有限完全交叉なので局所指標は $(1+u+u^2)^2$ となる。
このチャートの閉点は最低点だけであり、その長さが9である。反転すれば最高点も同じとなる。
残りの閉点はグラフの一意性と固定性から中央どうしである。
中央二直線の積には重み2の二関係だけが制限され、有限なので交叉長は2である。
既に全長を使うため、これはスキームとして残余全部である。
\end{proof}
